% !PRINTED 103
\weberpage{103}
\begin{center}
\textbf{§ 31. \gesperrt{Die linearen Fundamentaltransformationen
der $\vartheta$-Funktionen.}}
\end{center}

Bei der Anwendung auf die Transformation der $\Theta$-Funktionen
[§ 27, (12)] kommt zunächst der transformierte Modul
\[
\omega' = \dfrac{c + \partial\omega}{a + b\omega}
\]
in Betracht. Dieser ändert sich nicht, wenn die vier Transformationszahlen
 einen gemeinsamen Faktor $m$ bekommen. Die
Transformation heißt eine \gesperrt{eigentliche}, wenn $a, b, c, \partial$ keinen
gemeinsamen Faktor haben.

Das transformierte Argument
\[
u' = \dfrac{nu}{a + b\omega}
\]
geht über in $mu'$, wenn $a, b, c, \partial$ durch $ma, mb, mc, m\partial$, also $n$
durch $m^2n$ ersetzt wird. Die Multiplikation $\begin{pmatrix}m, & 0\\0, & m\end{pmatrix}$ läßt den
Modul $\omega$ ungeändert und verwandelt $u$ in $mu$.

Nach den Resultaten der beiden vorigen Paragraphen läßt
sich das ganze System der eigentlichen Transformationen herleiten
 durch wiederholte Anwendung der drei Transformationen
\[
\begin{pmatrix}1, & 0\\1, & 1\end{pmatrix}, \quad \begin{pmatrix}0, & 1\\-1, & 0\end{pmatrix}, \quad \begin{pmatrix}n, & 0\\0, & 1\end{pmatrix},
\]
und wir betrachten also zunächst die linearen Fundamentaltransformationen
 der $\vartheta$-Funktionen.

I. $\begin{pmatrix}1, & 0\\1, & 1\end{pmatrix}$ oder $(\omega, \omega + 1)$.

Nach § 27, (11), (12) ist
\weq{1}{\vartheta_{11}(u, \omega + 1) = A\vartheta_{11}(u, \omega),}
worin $A$ von $u$ unabhängig ist. Ersetzt man $u$ durch
\[
u + \tfrac{1}{2}, \quad u + \tfrac{\omega}{2}, \quad u + \tfrac{1+\omega}{2},
\]
so ergeben die Formeln § 21, (8)
\weq{2}{\vartheta_{10}(u, \omega + 1) = A\vartheta_{10}(u, \omega)}
\weq{3}{e^{\frac{\pi i}{4}}\vartheta_{00}(u, \omega + 1) = A\vartheta_{01}(u, \omega)}
\weq{4}{e^{\frac{\pi i}{4}}\vartheta_{01}(u, \omega + 1) = A\vartheta_{00}(u, \omega).}
